% Related work (0.5 p). Refs: SATNAC list (author-supplied) + web-verified (refs.bib verify_note)
\section{Background and Related Work}
\label{sec:related}

\textit{Aerial base stations and their testbeds.} The elevation-dependent
line-of-sight model~\cite{alhourani2014modeling} underlies most work on where
to place a UAV base station, from 3-D placement~\cite{boryaliniz2016efficient,
kalantari2016number} and trajectory design~\cite{wu2018common} to deep
reinforcement learning~\cite{qiu2020placement}; surveys cover the wider
field~\cite{zeng2016wireless,mozaffari2019tutorial}. These works optimise
coverage or rate under an analytic or simulated channel. Testbeds now run
OpenAirInterface on UAVs with O-RAN control~\cite{moore2024aerpaw} or emulate
5G mobility with ray tracing~\cite{iye2025owdt}. Our earlier
study~\cite{makhubela2026satnac} coupled a real 5G stack to a terrain
simulator and found that terrain decides whether moving the node helps.

\textit{Learning-based and autonomous management.} Zero-touch management aims
at fully autonomous operation under human oversight~\cite{coronado2022zsm}, and
O-RAN exposes closed control loops at sub-second to minute time scales in which
learned controllers can run~\cite{polese2023oran}; our three loops sit at the
same three scales, but outside the RAN.
Machine learning for link-quality estimation~\cite{cerar2021lqe}, fault
management~\cite{mulvey2019fault} and failure prediction in 5G RANs, with deep
and transformer models~\cite{islam2023dnn,hasan2024gentrap}, learns from the
history of a deployed network. A UAV arriving at a disaster site has no such
history for its terrain, so we test on an unseen terrain and in closed loop.

\textit{Buffering in cellular links.} Buffers that hold rather than drop
packets inflate delay~\cite{gettys2012bufferbloat}, and per-user cellular
buffers are a known source~\cite{jiang2012tackling}. Remedies range from active
queue management~\cite{nichols2012codel} and delay-aware rate
control~\cite{winstein2013sprout} to, recently, RAN-side congestion marking:
in 5G, time spent in the radio link control (RLC) buffer dominates one-way
delay~\cite{wan2025l4span}, retransmission and uplink-scheduling buffers can
add delay comparable to congestion~\cite{liu2025fog}, and sizing the RLC buffer
dynamically cuts that delay~\cite{irazabal2022buffer}. Those remedies change
network queues, the RAN or the transport protocol. Our cap applies the classic additive-increase,
multiplicative-decrease (AIMD) rule~\cite{chiu1989aimd} to a delay signal at the robot, because constant-rate
robot video does not back off and the UE's buffer is not ours to change.

% Draft table: what is new over SATNAC 2026 (sources: published SATNAC text,
% pasted by the author 2026-11-12; this paper's Sections III-V).
\begin{table}[t]
\caption{This Paper Versus Our Earlier Study~\cite{makhubela2026satnac}}
\label{tab:satnac}
\centering\footnotesize\setlength{\tabcolsep}{3pt}
\begin{tabular}{@{}p{0.2\columnwidth}p{0.36\columnwidth}p{0.38\columnwidth}@{}}
\toprule
 & SATNAC 2026 & This paper \\
\midrule
Robots & 2 & 4 \\
UAV control & fixed hover vs.\ random walk & forecaster, planner, altitude learners vs.\ hover and centroid \\
Management loops & rule-based coverage guard & uplink cap, liveness monitor, learned x--y--z placement \\
Traffic & 2-s 8\,Mb/s bursts per cycle & continuous 1.5\,Mb/s video per robot \\
C2 available & $\geq$1 of 5 echoes answered & ping loss $<$\,50\,\% \\
Validation & static vs.\ roaming phase, one run per terrain & held-out terrain, null controls, repeated seeds \\
Main finding & mobility helps in open field, hurts in forest & bound queues and check liveness first; learning ties simple rules \\
\bottomrule
\end{tabular}
\end{table}

